\chapter{Literature Review}\label{ch:literature}

This chapter reviews the work most relevant to the thesis questions: can a
system detect and name ECG abnormalities, and how does performance change when
earlier ECG from the same patient is available? The review focuses on four
areas: ECG datasets and evaluation protocols, supervised heartbeat
classification, normal-only anomaly detection, and statistical validation. It
is a focused narrative review, not a systematic review or meta-analysis. Papers
were selected because they explain a design choice, show an important
alternative, or reveal a risk in patient-independent evaluation.

Headline accuracy is not used to rank the papers. ECG studies often use
different leads, labels, record lengths, patient splits, and output units. For
example, a patient-specific beat classifier, a 12-lead diagnostic model, and
the six-class system in this thesis answer different questions. Their accuracy
values cannot be compared directly. This matters because a high score may come
from an easier class mapping, or from placing beats from the same person in
both training and test data.

PhysioNet provides open physiological databases and common software tools
\cite{goldberger2000physionet}. The MIT--BIH Arrhythmia Database (MITDB) became
a widely used benchmark because it contains expert beat annotations and two ECG
channels \cite{moody2001impact,mitdb}. The MIT--BIH Normal Sinus Rhythm
Database provides 18 real long-term recordings from subjects without
significant arrhythmias \cite{nsrdb}. It is suitable for normal-only learning,
but its headers identify the channels only as ECG1 and ECG2. The present study
therefore uses one fixed first channel and reports this uncertainty instead of
assigning an unsupported anatomical lead name.

The evaluation protocol is as important as the database. A beat-wise random
split can place very similar beats from one patient in both training and test
sets. An inter-patient split instead tests whether the learned boundary transfers
to a different person. De Chazal et al. established an influential inter-patient
heartbeat study using ECG morphology and heartbeat-interval features
\cite{dechazal2004heartbeat}. Later research continued to treat patient variation
as a central problem rather than assuming that all beats are independent
\cite{mousavi2019seq2seq,zhang2021interpatient}.

This evidence supports a clearly defined evaluation protocol instead of a
single ambiguous score. The final experiment keeps chronological order, trains
on the first 80\% of every record, and tests on the final 20\%. This protocol
is stricter than a random beat split. However, it measures patient-adapted
temporal continuation, not performance on a new patient. The separate NSRDB
cohort is used only for the normal autoencoder; no healthy record trains either
supervised tree model.

Early automatic ECG systems used measured morphology and timing. De Chazal et
al. combined waveform information with RR intervals, showing why beat shape and
beat timing should be treated as complementary evidence
\cite{dechazal2004heartbeat}. This idea remains useful even when the classifier
changes: morphology helps describe PVC and bundle-branch beats, while RR history
helps describe premature beats and irregular rhythms.

Huang et al. studied inter-patient separation of normal, LBBB, and RBBB beats
with ensemble classifiers \cite{huang2014bbb}. That study is particularly
relevant because bundle-branch morphology can persist through a record. However,
its three-class task is narrower than the present N/PVC/PAC/LBBB/RBBB/AFib task.
Its results therefore justify using morphology and ensemble learning, but they
cannot be copied as an expected accuracy for this thesis.

Feature-based models have two practical advantages for this dataset. They make
the evidence easier to inspect, and they can be controlled when only a few
patients contain a class. Their weakness is that hand-designed summaries may
miss useful waveform detail. This trade-off led to the combined design used
here: pooled waveform and RR features are kept, while an autoencoder adds
learned normal-reconstruction evidence.

One-dimensional convolutional neural networks can learn ECG morphology directly.
Kiranyaz et al. developed a patient-specific real-time classifier
\cite{kiranyaz2016ecgcnn}. Its patient adaptation is useful when labelled beats
from the target person exist. It is therefore directly relevant to the
final patient-adapted setting, but it does not answer performance before
seeing a new patient. Acharya et al. classified arrhythmias
from ECG segments with a convolutional network \cite{acharya2017cnn}. Segment-
level work demonstrates the value of learned temporal features, but its label
definition and split must match the target study before accuracy can be compared.

Large deep-learning studies have also reported strong performance on ambulatory
single-lead recordings and 12-lead diagnostic ECGs
\cite{hannun2019cardiologist,ribeiro2020ecg}. These works show the potential of
deep networks when many patients and labels are available. They also show why
dataset scale matters. Their databases, lead systems, and diagnostic outputs are
different from the small six-class MITDB experiment, so they are evidence for a
method family rather than a performance baseline for this thesis.

Sequence models use neighbouring beats to add rhythm context. Mousavi et al.
compared intra-patient and inter-patient heartbeat classification with a
sequence-to-sequence approach \cite{mousavi2019seq2seq}. Zhang et al. added an
adversarial subject classifier so that the learned representation would contain
less patient identity \cite{zhang2021interpatient}. These approaches address
generalisation directly, but they also need enough independent patients. The
present study has only 45 eligible subject clusters, so the final system uses a
simpler tree hierarchy. This is a data-availability decision, not a claim that
deep learning is generally inferior.

Autoencoders learn a compressed representation by reconstructing their input
\cite{hinton2006autoencoder}. Denoising training corrupts the input but retains a
clean target, encouraging the network to learn stable structure rather than an
exact copy \cite{vincent2008denoising}. In anomaly detection, the model is often
trained on normal sequences and reconstruction error is used as evidence that a
new sequence differs from the learned normal pattern
\cite{chandola2009anomaly,malhotra2016lstm}.

Temporal convolutional autoencoders use convolution and dilation to cover longer
time patterns efficiently. Thill et al. evaluated this idea on time-series
anomaly tasks that included ECG \cite{thill2021tcnae}. ECG-specific studies have
also investigated convolutional variational and adversarial autoencoders
\cite{jang2021cvae,zhou2022aae}. Together, these studies support using a
convolutional autoencoder for normal ECG representation.

The literature does not mean that reconstruction error is a diagnosis. A
high-capacity model may reconstruct an abnormal input well, while noise or a
different electrode position may increase the error of a normal input. For that
reason, this thesis does not use one reconstruction threshold as the final
answer. It uses four residual summaries as features in a supervised abnormal
detector. The normal-only idea is kept, but the final decision can combine
residual, morphology, spectral, and RR evidence.

ECG datasets contain unequal class counts and unequal numbers of patients per
class. A high overall accuracy can therefore hide failure on a rare arrhythmia.
Precision--recall analysis is more informative than ROC analysis in strongly
imbalanced binary problems \cite{saito2015pr}, while class-balanced learning can
increase the influence of rare classes \cite{cui2019classbalanced}. This thesis
therefore reports sensitivity, specificity, balanced accuracy, macro F1, and
per-class recall in addition to accuracy.

Thousands of beats from a small number of subjects do not become thousands of
independent patients. If individual beats are resampled, the confidence
intervals become too narrow. Patient-cluster bootstrap methods keep
within-patient dependence \cite{rutter2000clustered}. The final experiment gives
each of 45 subject clusters equal weight and resamples complete subjects. This
also follows modern guidance to report the data split, model development, and
validation population clearly \cite{collins2024tripodai}.

The literature provides strong individual ideas: normal-only reconstruction,
learned waveform features, morphology and RR features, patient-invariant
learning, specialised bundle-branch classification, and clustered statistical
validation. However, reported results usually address only part of the present
problem. They may use a different class mapping, adapt to the test patient, use
more leads, or evaluate anomaly detection without naming the abnormal type.

The gap addressed by this thesis is therefore not simply the need for another
classifier. The need is to evaluate a complete two-stage system while making
the generalisation claim explicit:
\begin{enumerate}
    \item learn normal reconstruction from one fixed channel of all 18 real
    MIT--BIH normal-sinus recordings;
    \item use a chronological 80/20 experiment for patient-adapted temporal evidence;
    \item keep each final test portion unavailable during selection; and
    \item report separate abnormal, subtype, whole-system, patient-level, and
    uncertainty results.
\end{enumerate}

Table~\ref{tab:literature-comparison} summarises the papers that contributed most
directly to the data protocol, model architecture, and evaluation method used in
this research.

\begingroup
\footnotesize
\setlength{\tabcolsep}{4pt}
\begin{longtable}{>{\raggedright\arraybackslash}p{5.0cm}>{\raggedright\arraybackslash}p{3.4cm}>{\raggedright\arraybackslash}p{5.3cm}}
\caption{Main papers referenced in this study and their contribution.}\label{tab:literature-comparison}\\
\toprule
Paper title & Author(s) & Main contribution used in this research\\
\midrule
\endfirsthead
\toprule
Paper title & Author(s) & Main contribution used in this research\\
\midrule
\endhead
MIT--BIH Normal Sinus Rhythm Database \cite{nsrdb} & George B. Moody & Provides the 18 real long-term recordings without significant arrhythmias used for normal-only autoencoder development; its generic channel names also motivate explicit reporting of lead uncertainty.\\
Automatic Classification of Heartbeats Using ECG Morphology and Heartbeat Interval Features \cite{dechazal2004heartbeat} & Philip de Chazal, Maria O'Dwyer, and Richard B. Reilly & Established the value of combining ECG morphology with RR intervals and carefully defining the evaluation population. This supports the morphology, timing, and reporting choices of the present system.\\
Reducing the Dimensionality of Data with Neural Networks \cite{hinton2006autoencoder} & Geoffrey E. Hinton and Ruslan R. Salakhutdinov & Provides the representation-learning basis of the autoencoder used to learn normal ECG structure.\\
Extracting and Composing Robust Features with Denoising Autoencoders \cite{vincent2008denoising} & Pascal Vincent, Hugo Larochelle, Yoshua Bengio, and Pierre-Antoine Manzagol & Supports adding noise to the input while keeping a clean reconstruction target, so the autoencoder learns stable normal patterns.\\
Temporal Convolutional Autoencoder for Unsupervised Anomaly Detection in Time Series \cite{thill2021tcnae} & Markus Thill, Wolfgang Konen, and Thomas Baeck & Supports temporal convolution, dilation, and reconstruction error for detecting unusual time-series patterns, including ECG signals.\\
Inter- and Intra-patient ECG Heartbeat Classification for Arrhythmia Detection: A Sequence to Sequence Deep Learning Approach \cite{mousavi2019seq2seq} & Sajad Mousavi and Fatemeh Afghah & Shows the importance of neighbouring-beat rhythm context and clearly separates intra-patient from inter-patient evaluation.\\
Interpatient ECG Heartbeat Classification with an Adversarial Convolutional Neural Network \cite{zhang2021interpatient} & Jing Zhang, Aiping Liu, Deng Liang, Xun Chen, and Min Gao & Explains patient-identity leakage and why results from known-patient temporal tests must not be presented as new-patient performance.\\
Detection of Inter-Patient Left and Right Bundle Branch Block Heartbeats in ECG Using Ensemble Classifiers \cite{huang2014bbb} & Huifang Huang, Jie Liu, Qiang Zhu, Rui Wang, and Guangshu Hu & Supports morphology and ensemble classifiers for LBBB and RBBB, while showing that this specialised three-class task is narrower than the thesis task.\\
The Precision--Recall Plot Is More Informative than the ROC Plot When Evaluating Binary Classifiers on Imbalanced Datasets \cite{saito2015pr} & Takaya Saito and Marc Rehmsmeier & Justifies reporting precision, recall, and F1 because the ECG classes are imbalanced and accuracy alone can hide failure.\\
Bootstrap Estimation of Diagnostic Accuracy with Patient-Clustered Data \cite{rutter2000clustered} & Carolyn M. Rutter & Provides the reason for resampling complete patient records instead of treating correlated beats as independent observations.\\
TRIPOD+AI Statement: Updated Guidance for Reporting Clinical Prediction Models That Use Regression or Machine Learning Methods \cite{collins2024tripodai} & Gary S. Collins et al. & Supports transparent reporting of the patient split, model-development process, temporal evaluation, limitations, and intended use.\\
\bottomrule
\end{longtable}
\endgroup
